\subsection{拓扑除环与域}

在下文以及第四、第五章中，若 K 是一个除环，我们用 \(K^{*}\) 表示 K 的非零元素所成的乘法群。

定义 4. 拓扑除环是一个集合 K，其上带有除环结构与一个和 K 的环结构相容的拓扑，且此外还满足下列公理：

(KT) 映射 \(x \to x^{-1}\)（\(\mathbf{K}^*\) 到 \(\mathbf{K}^*\) 中的映射）是连续的。

交换的拓扑除环称为拓扑域。

集合 K 上的除环结构与一个拓扑称为相容的，如果相应的环结构与该拓扑相容，且此外公理 (KT) 成立。

例. 1) 在任一除环 K 上，离散拓扑与除环结构相容。拓扑为离散的拓扑除环称为离散除环。

* 2) 有理直线 Q（相应地，实直线 R）的拓扑与 R（相应地，R）的域结构相容（见第四章，§ 3）。*

定义 4 表明，若 K 是拓扑除环，则由 K 的拓扑在乘法群 \(K^{*}\) 上诱导的拓扑与 \(K^{*}\) 的群结构相容。

若 \(a \neq 0\)，则位似 \(x \rightarrow ax\) 与 \(x \rightarrow xa\) 都是 K 到其自身上的同胚；且映射 \(x \rightarrow ax + b\) 对所有 \(b \in K\) 亦然。注意，位似 \(x \rightarrow ax\) 与 \(x \rightarrow xa\) 在 \(a \neq 0\) 时是 K 的（拓扑）加法群的自同构。由此，若 V 是 K 中 o 的任一邻域，则对所有 \(a \neq 0\)，aV 与 Va 都是 o 的邻域。

设 X 是拓扑空间，f 是 X 到拓扑除环 K 中的一个映射。若 f 在一点 \(x_{0} \in X\) 连续且 \(f(x_{0}) \neq 0\)，则 \(f^{-1}\) 在 \(x_{0}\) 连续。特别地，若 K 是拓扑域，则每个 n 个变元、系数在 K 中的有理函数，都在 \(K^{n}\) 上其分母不为零的每一点连续。

同样，若 \(f\) 是一个集合 \(\mathbf{X}\)（由滤子 \(\mathfrak{F}\) 所滤化）到豪斯多夫拓扑除环 \(\mathbf{K}\) 中的映射，且 \(\lim_{\mathfrak{F}} f\) 存在且不为 \(o\)，则 \(\lim_{\mathfrak{F}} f^{-1}\) 存在，并且我们有

\[
\lim _ {\mathfrak {F}} f ^ {- 1} = (\lim _ {\mathfrak {F}} f) ^ {- 1}.\tag{4}
\]

若 H 是拓扑除环 K 的除子环，则由 K 的拓扑在 H 上诱导的拓扑与 H 的除环结构相容。这样在 H 上定义的拓扑除环结构称为由 K 的拓扑除环结构所诱导。此外，\(\overline{H}\) 也是 K 的除子环（证明与命题 5 的证明类似）。

在拓扑除环 K 中，集合 \(\{0\}\) 的闭包由命题 5 是一个双边理想，因而必须是 \(\{0\}\) 或 K。换句话说，若 K 的拓扑不是最粗拓扑（第一章，§ 2，第 2 节），则它是豪斯多夫的（§ 1，第 2 节，命题 2）。

